\subsection{The full nonconstant unit-fibre calculation: profinite orbit
coordinates, connecting maps, and the natural middle algebra}
\label{subsec:nonconstant-unit-fibre-k-theory}

The stable extension embedding of the preceding subsection sees the constant
subgroup
\(\mathbb Z\mathbf 1_{K_P}\subset C(K_P,\mathbb Z)\).  We now calculate the
remaining classes rather than referring to them collectively as
``nonconstant.''  The answer has three distinct pieces: invariant functions
on finite-place profinite quotients, coinvariants with nontrivial finite-level
transition multiplicities, and an archimedean sign-even subgroup.  All three
are needed for the K-theory of the natural one-zero algebra.

Fix a nonempty finite prime set \(P\).  For \(p\in P\), put
\begin{equation}
 Y_{p,P}=K_{P\setminus p}
 =\prod_{q\in P\setminus\{p\}}\mathbb Z_q^\times,
 \qquad
 g_{p,P}=(p)_{q\ne p}\in Y_{p,P}.
 \label{eq:unit-fibre-Y-and-g}
\end{equation}
The empty product is one point.  Let
\begin{equation}
 a_{p,P}(y)=g_{p,P}y,
 \qquad
 (\sigma_{p,P}f)(y)=f(g_{p,P}^{-1}y),
 \label{eq:unit-fibre-translation-automorphism}
\end{equation}
and let
\begin{equation}
 H_{p,P}=\overline{\langle g_{p,P}\rangle}\subset Y_{p,P},
 \qquad
 q_{p,P}:Y_{p,P}\longrightarrow Y_{p,P}/H_{p,P}
 \label{eq:unit-fibre-procyclic-closure}
\end{equation}
be the closed procyclic subgroup and quotient map.  If \(Y_{p,P}\) is not a
point, the \(\mathbb Z\)-action generated by \(a_{p,P}\) is free: an equality
\(p^ny=y\) in one factor \(\mathbb Z_q^\times\) would imply
\(p^n=1\) in \(\mathbb Q_q\), and hence as a rational number, so \(n=0\).
For \(P=\{p\}\), the action on the one-point \(Y_{p,P}\) is instead trivial;
this exceptional stabilizer is retained below and gives the ordinary circle.

\begin{lemma}[Mapping-torus model and complete finite-place K-groups]
\label{lem:unit-fibre-mapping-torus-k-groups}
There is a literal isomorphism
\begin{equation}
 C(B_{p,P})\cong M_{\sigma_{p,P}}
 =\{F\in C([0,1],C(Y_{p,P})):F(1)=\sigma_{p,P}(F(0))\}.
 \label{eq:unit-fibre-mapping-torus-isomorphism}
\end{equation}
Under this isomorphism and the Green-module coordinate for
\(\mathfrak D_{p,P}\),
\begin{align}
 K_0(\mathfrak D_{p,P})
 &\cong
 \ker(1-\sigma_{p,P,*})
 \cong C(Y_{p,P}/H_{p,P},\mathbb Z),
 \label{eq:finite-edge-K0-invariants}\\
 K_1(\mathfrak D_{p,P})
 &\cong
 \operatorname{coker}(1-\sigma_{p,P,*})
 \cong
 \frac{C(Y_{p,P},\mathbb Z)}
 {(1-\sigma_{p,P})C(Y_{p,P},\mathbb Z)}.
 \label{eq:finite-edge-K1-coinvariants}
\end{align}
No constant-function reduction occurs in either group.
\end{lemma}

\begin{proof}
Write \(\ell_p=\log p\).  For \(G\in C(B_{p,P})\), define
\[
 (\Theta G)(u)(y)=G([y,u\ell_p]),\qquad 0\leq u\leq1.
\]
Since \([y,\ell_p]=[g_{p,P}^{-1}y,0]\),
\[
 (\Theta G)(1)(y)=(\Theta G)(0)(g_{p,P}^{-1}y)
 =\sigma_{p,P}((\Theta G)(0))(y).
\]
Conversely, the endpoint equation makes a path in the displayed mapping
torus descend uniquely to the quotient
\((Y_{p,P}\times\mathbb R)/\langle\beta_p\rangle\).  These two maps preserve
multiplication, involution, and the supremum norm, proving
\eqref{eq:unit-fibre-mapping-torus-isomorphism}.

We also record the coefficient groups rather than hiding them behind a
topological K-theory identification.  For every compact profinite space
\(Y\), a projection over \(Y\) is, after refining a finite clopen partition,
unitarily equivalent to a constant matrix projection on each cell.  Stable
equivalence is therefore determined exactly by its locally constant rank
function, so
\[
 K_0(C(Y))=C(Y,\mathbb Z).
\]
Moreover every continuous unitary matrix over \(Y\) is uniformly approximated
within distance less than two by one that is constant on a finite clopen
partition.  The two unitaries are homotopic, and each constant value is
homotopic to the identity inside a unitary group.  Thus
\(K_1(C(Y))=0\).

Blackadar's mapping-torus exact sequence and its explicit connecting maps are
\(1-\sigma_*\) under the Bott coordinate
\cite[Theorem~10.2.1, Definition~10.3.1, and Proposition~10.4.1,
pp.~83--87]{blackadar1986}.  Substitution of the two coefficient groups just
proved gives the kernel and cokernel in
\eqref{eq:finite-edge-K0-invariants}--
\eqref{eq:finite-edge-K1-coinvariants}.  A locally constant function is fixed
by \(\sigma_{p,P}\) exactly when it is constant on every orbit of
\(g_{p,P}\), hence on every coset of the closure \(H_{p,P}\).  Pullback along
\(q_{p,P}\) therefore identifies the kernel with
\(C(Y_{p,P}/H_{p,P},\mathbb Z)\).  Finally, the restricted Green module from
Theorem~\ref{thm:multiplace-finite-edge-coordinates} transports these groups
to \(\mathfrak D_{p,P}\) without changing the extension boundary coordinate.
\end{proof}

The coinvariant group in \eqref{eq:finite-edge-K1-coinvariants} generally
does not reduce to a free group on a final orbit set.  Its exact finite-level
coordinates are as follows.  Let \(\mathcal M_{p,P}\) be the directed set of
tuples \(\mathbf m=(m_q)_{q\ne p}\), with \(m_q\geq1\), ordered
coordinatewise.  For an empty tuple set all objects below equal to the
one-point object and put \(L_{\varnothing}=1\).  Otherwise define
\begin{equation}
 G_{\mathbf m}=\prod_{q\ne p}
 (\mathbb Z/q^{m_q}\mathbb Z)^\times,
 \quad
 g_{\mathbf m}=(p\bmod q^{m_q})_{q\ne p},
 \quad
 L_{\mathbf m}=\operatorname{ord}_{G_{\mathbf m}}(g_{\mathbf m}),
 \quad
 \Omega_{\mathbf m}=G_{\mathbf m}/\langle g_{\mathbf m}\rangle.
 \label{eq:finite-orbit-coordinate-data}
\end{equation}
Every orbit has exactly \(L_{\mathbf m}\) points.  Thus the stabilizer of a
point for the finite-level \(\mathbb Z\)-action is precisely
\(L_{\mathbf m}\mathbb Z\); these finite-level stabilizers disappear in the
nontrivial inverse limit but are not deleted from the transition maps.

For an orbit \(O\in\Omega_{\mathbf m}\), let
\[
 u_O=\mathbf1_O
 \quad\hbox{in the invariant group},
 \qquad
 v_O=[\delta_x]\quad(x\in O)
 \quad\hbox{in the coinvariant group}.
\]
The second class is independent of \(x\in O\).  If
\(\mathbf m'\geq\mathbf m\), let
\(\pi_{\mathbf m',\mathbf m}:G_{\mathbf m'}\to G_{\mathbf m}\) be reduction.
Then \(L_{\mathbf m}\mid L_{\mathbf m'}\), and the two transition maps are
\begin{align}
 \iota^0_{\mathbf m',\mathbf m}(u_O)
 &=\sum_{\substack{O'\in\Omega_{\mathbf m'}\\
                    O'\subset\pi_{\mathbf m',\mathbf m}^{-1}(O)}}u_{O'},
 \label{eq:finite-invariant-transition}\\
 \iota^1_{\mathbf m',\mathbf m}(v_O)
 &=\frac{L_{\mathbf m'}}{L_{\mathbf m}}
   \sum_{\substack{O'\in\Omega_{\mathbf m'}\\
                    O'\subset\pi_{\mathbf m',\mathbf m}^{-1}(O)}}v_{O'}.
 \label{eq:finite-coinvariant-transition}
\end{align}

\begin{proposition}[Finite-level coordinate form, including multiplicities]
\label{prop:unit-fibre-finite-level-limits}
The two groups in
\eqref{eq:finite-edge-K0-invariants}--
\eqref{eq:finite-edge-K1-coinvariants} are respectively
\begin{align}
 C(Y_{p,P}/H_{p,P},\mathbb Z)
 &\cong
 \varinjlim_{\mathbf m}
 (\mathbb Z[\Omega_{\mathbf m}],\iota^0),
 \label{eq:finite-invariant-direct-limit}\\
 C(Y_{p,P},\mathbb Z)_{\sigma_{p,P}}
 &\cong
 \varinjlim_{\mathbf m}
 (\mathbb Z[\Omega_{\mathbf m}],\iota^1).
 \label{eq:finite-coinvariant-direct-limit}
\end{align}
Both displayed transition maps are injective.  In particular the
coinvariant group is torsion-free, but no free decomposition or canonical
splitting is asserted.
\end{proposition}

\begin{proof}
Every locally constant integer-valued function on \(Y_{p,P}\) factors through
some \(G_{\mathbf m}\), and pullback along the reduction maps gives
\[
 C(Y_{p,P},\mathbb Z)
 =\varinjlim_{\mathbf m}\mathbb Z^{G_{\mathbf m}}.
\]
At level \(\mathbf m\), the invariant functions have basis
\((u_O)_{O\in\Omega_{\mathbf m}}\), and pullback of \(\mathbf1_O\) is the
sum of the characteristic functions of all fine orbits over \(O\).  This is
\eqref{eq:finite-invariant-transition}.

The coinvariants of the permutation module of one orbit are \(\mathbb Z\),
with every point mass representing the same generator \(v_O\).  Fix
\(x\in O\).  Pullback of \(\delta_x\) is the sum of the point masses over
\(\pi^{-1}(x)\).  Every fine orbit \(O'\) above \(O\) maps onto \(O\) with
degree \(L_{\mathbf m'}/L_{\mathbf m}\), so exactly that many of its points
lie above \(x\).  Hence its contribution is
\((L_{\mathbf m'}/L_{\mathbf m})v_{O'}\), proving
\eqref{eq:finite-coinvariant-transition}.  Equivalently, the number of fine
orbits above \(O\) is
\[
 |\ker\pi_{\mathbf m',\mathbf m}|
 \frac{L_{\mathbf m}}{L_{\mathbf m'}},
\]
which verifies that all points in the pullback have been counted exactly
once.

Filtered colimits commute with the cokernel of
\(1-\sigma_{p,P}\), giving
\eqref{eq:finite-coinvariant-direct-limit}.  For invariants, a function that
factors through \(G_{\mathbf m}\) is invariant in the inverse limit exactly
when its finite-level function is constant on the
\(\langle g_{\mathbf m}\rangle\)-orbits, which gives
\eqref{eq:finite-invariant-direct-limit}.  Distinct coarse orbits have
disjoint inverse images.  Both transition formulas therefore send a nonzero
integer combination to a nonzero combination on disjoint supports.  They are
injective.  An injective filtered union of torsion-free groups is
torsion-free, proving the final statement.
\end{proof}

We next compute the connecting map on every invariant rank function.  Let
\begin{equation}
 \operatorname{pr}_{p,P}:K_P\longrightarrow Y_{p,P}
 \label{eq:finite-unit-drop-map}
\end{equation}
drop the \(p\)-coordinate, and define
\begin{equation}
 \iota_{p,P}:C(Y_{p,P}/H_{p,P},\mathbb Z)
 \longrightarrow C(K_P,\mathbb Z),
 \qquad
 \iota_{p,P}(h)=h\circ q_{p,P}\circ\operatorname{pr}_{p,P}.
 \label{eq:finite-invariant-ideal-pullback}
\end{equation}

\begin{theorem}[Full finite-edge connecting map and middle K-groups]
\label{thm:full-finite-edge-boundary-map}
In the coordinates above, the connecting map of the \(p\)-edge is
\begin{equation}
 \partial_{p,P}(h)=-\iota_{p,P}(h)
 \qquad
 (h\in C(Y_{p,P}/H_{p,P},\mathbb Z)).
 \label{eq:full-finite-edge-boundary-map}
\end{equation}
It is injective.  Consequently
\begin{equation}
 K_0(A_{\varnothing p}^{(P)})=0
 \label{eq:finite-edge-middle-K0-zero}
\end{equation}
and there is an exact sequence
\begin{equation}
 0\longrightarrow
 \frac{C(K_P,\mathbb Z)}
      {\iota_{p,P}C(Y_{p,P}/H_{p,P},\mathbb Z)}
 \longrightarrow K_1(A_{\varnothing p}^{(P)})
 \longrightarrow C(Y_{p,P},\mathbb Z)_{\sigma_{p,P}}
 \longrightarrow0.
 \label{eq:finite-edge-middle-K1-exact}
\end{equation}
No splitting of \eqref{eq:finite-edge-middle-K1-exact} is selected.
\end{theorem}

\begin{proof}
It is enough by additivity to take
\(h=\mathbf1_E\) for a clopen subset
\(E\subset Y_{p,P}/H_{p,P}\).  Put
\(\widetilde E=q_{p,P}^{-1}(E)\).  It is clopen and invariant under
\(g_{p,P}\).  The corresponding clopen part of the mapping-torus boundary is
\[
 B_E=(\widetilde E\times\mathbb R)/\langle\beta_p\rangle
 \subset B_{p,P}.
\]
In the orbit coordinates of
Theorem~\ref{thm:multiplace-finite-edge-coordinates}, the scalar function
\begin{equation}
 x_E(a,t)=
 \frac{\mathbf1_{\widetilde E}(\operatorname{pr}_{p,P}(a))}{1+e^t},
 \qquad
 x_E([b,s])=\mathbf1_{\widetilde E}(b)
 \label{eq:finite-edge-explicit-lift}
\end{equation}
is continuous.  Invariance of \(\widetilde E\) makes the boundary formula
independent of the representative \((b,s)\), while the exact cross-convergence
criterion makes the generic value tend to it as \(t\to-\infty\).  The only
escaping end is \(t\to+\infty\), where the function vanishes.  Thus \(x_E\)
is a self-adjoint lift of the boundary projection \(\mathbf1_{B_E}\).

Its exponential is the identity off
\(\operatorname{pr}_{p,P}^{-1}(\widetilde E)\).  On every retained generic
unit fibre above \(\widetilde E\), its argument decreases from \(2\pi\) to
zero as \(t\) increases, so it has winding \(-1\) in the previously declared
orientation.  Hence
\[
 \partial_{p,P}[\mathbf1_{B_E}]
 =-\mathbf1_{\operatorname{pr}_{p,P}^{-1}(\widetilde E)}.
\]
Every locally constant integer function is an integer combination of such
characteristic functions, proving
\eqref{eq:full-finite-edge-boundary-map}.  The map
\(q_{p,P}\circ\operatorname{pr}_{p,P}\) is onto, so its pullback is injective.

Finally, \(\mathfrak I_P\) is Morita equivalent to
\(C_0(K_P\times\mathbb R)\), whose K-groups in the retained suspension
coordinate are
\[
 K_0(\mathfrak I_P)=0,
 \qquad K_1(\mathfrak I_P)=C(K_P,\mathbb Z).
\]
Insert these, Lemma~\ref{lem:unit-fibre-mapping-torus-k-groups}, and the
injective connecting map into the edge six-term sequence.  Exactness gives
\eqref{eq:finite-edge-middle-K0-zero} and
\eqref{eq:finite-edge-middle-K1-exact}.
\end{proof}

The archimedean summand contributes all sign-even functions, not only its
constant unit.  Retain the clopen half \(H_P^+\subset K_P\) constructed in
the proof of Theorem~\ref{thm:multiplace-arch-boundary}, so
\begin{equation}
 K_P=H_P^+\sqcup(-H_P^+).
 \label{eq:arch-clopen-half}
\end{equation}
Define the continuous representative map
\begin{equation}
 r_P:K_P\longrightarrow H_P^+,
 \qquad
 r_P(c)=
 \begin{cases}c,&c\in H_P^+,\\-c,&c\in-H_P^+.
 \end{cases}
 \label{eq:arch-representative-map}
\end{equation}

\begin{theorem}[Full archimedean boundary map]
\label{thm:full-archimedean-unit-fibre-boundary}
The explicit matrix coordinate gives
\begin{equation}
 K_0(\mathfrak D_{\infty,P})\cong C(H_P^+,\mathbb Z),
 \qquad
 K_1(\mathfrak D_{\infty,P})=0.
 \label{eq:arch-boundary-full-K-groups}
\end{equation}
Under the generic coordinate
\(K_1(\mathfrak I_P)=C(K_P,\mathbb Z)\), the complete archimedean
connecting map is
\begin{equation}
 \partial_{\infty,P}(h)=-r_P^*h.
 \label{eq:full-archimedean-boundary-map}
\end{equation}
Its image is exactly the sign-even subgroup
\begin{equation}
 C(K_P,\mathbb Z)^{\mathrm{ev}}
 =\{f:f(c)=f(-c)\text{ for all }c\in K_P\}.
 \label{eq:sign-even-unit-functions}
\end{equation}
\end{theorem}

\begin{proof}
The clopen decomposition \eqref{eq:arch-clopen-half} trivializes the free
double cover, and the already displayed orbit matrices give
\[
 C(K_P)\rtimes C_2\cong M_2(C(H_P^+)).
\]
The profinite K-theory calculation in the proof of
Lemma~\ref{lem:unit-fibre-mapping-torus-k-groups} yields
\eqref{eq:arch-boundary-full-K-groups}.

For the sign and orientation, write every \(c\in K_P\) uniquely as
\(c=\varepsilon h\), with \(h\in H_P^+\) and
\(\varepsilon\in\{1,-1\}\).  The orbit map for the diagonal action on
\(K_P\times\mathbb R\) is explicitly
\begin{equation}
 [(\varepsilon h,z)]\longmapsto(h,w=\varepsilon z)
 \quad\hbox{in }H_P^+\times\mathbb R.
 \label{eq:arch-orbit-coordinate-w}
\end{equation}
On the generic part \(w\ne0\), its relation with the retained ideal coordinate
\((a,t)\in K_P\times\mathbb R\) is
\begin{equation}
 a=\operatorname{sgn}(w)h,
 \qquad t=\log|w|.
 \label{eq:arch-w-to-generic-coordinate}
\end{equation}

For a clopen \(E\subset H_P^+\), the function
\[
 x_E(h,w)=\frac{\mathbf1_E(h)}{1+|w|}
\]
on the commutative orbit extension is a self-adjoint lift of
\(\mathbf1_E\) at \(w=0\).  On both components \(w>0\) and \(w<0\), the
argument of its exponential decreases from \(2\pi\) to zero as
\(t=\log|w|\) increases.  Formula
\eqref{eq:arch-w-to-generic-coordinate} therefore gives pointwise winding
\(-1\) on
\(E\sqcup(-E)=r_P^{-1}(E)\) and zero elsewhere.  The extension-level Green
module transports this calculation to the crossed-product extension, so
\[
 \partial_{\infty,P}[\mathbf1_E]
 =-\mathbf1_{r_P^{-1}(E)}.
\]
Additivity proves \eqref{eq:full-archimedean-boundary-map}.  Pullback by
\(r_P\) is injective and its image consists exactly of the functions constant
on each pair \(\{c,-c\}\), proving
\eqref{eq:sign-even-unit-functions}.
\end{proof}

We can now assemble all places.  Abbreviate
\begin{equation}
 \mathcal F_{p,P}=C(Y_{p,P}/H_{p,P},\mathbb Z),
 \qquad
 \mathcal C_{p,P}=C(Y_{p,P},\mathbb Z)_{\sigma_{p,P}}.
 \label{eq:finite-invariant-coinvariant-abbreviations}
\end{equation}

\begin{theorem}[Full K-theory and connecting map of the natural one-zero
extension]
\label{thm:full-natural-one-zero-k-theory}
The quotient groups are
\begin{align}
 K_0(\mathfrak D_P)
 &\cong
 \left(\bigoplus_{p\in P}\mathcal F_{p,P}\right)
 \oplus C(H_P^+,\mathbb Z),
 \label{eq:global-natural-boundary-K0}\\
 K_1(\mathfrak D_P)
 &\cong\bigoplus_{p\in P}\mathcal C_{p,P}.
 \label{eq:global-natural-boundary-K1}
\end{align}
In these exact coordinates the connecting map is
\begin{equation}
 \partial_P^{(1)}((h_p)_{p\in P},h_\infty)
 =-\sum_{p\in P}\iota_{p,P}(h_p)-r_P^*h_\infty
 \quad\text{in }C(K_P,\mathbb Z).
 \label{eq:global-full-unit-fibre-boundary-map}
\end{equation}

For a finite tuple \(h=(h_p)_{p\in P}\), put
\begin{equation}
 S_h(c)=\sum_{p\in P}\iota_{p,P}(h_p)(c).
 \label{eq:finite-place-sum-function}
\end{equation}
Then projection to the finite-place coordinates gives an isomorphism
\begin{equation}
 K_0(\mathfrak N_P)
 \xrightarrow{\ \cong\ }
 \left\{h\in\bigoplus_{p\in P}\mathcal F_{p,P}:
 S_h(c)=S_h(-c)\text{ for every }c\in K_P\right\},
 \label{eq:natural-middle-K0-exact-kernel}
\end{equation}
whose inverse is
\begin{equation}
 h\longmapsto((h_p)_p,-S_h|_{H_P^+}).
 \label{eq:natural-middle-K0-kernel-inverse}
\end{equation}

Define
\begin{equation}
 d_{p,P}:\mathcal F_{p,P}\longrightarrow C(H_P^+,\mathbb Z),
 \qquad
 d_{p,P}(h)(c)=\iota_{p,P}(h)(c)-\iota_{p,P}(h)(-c),
 \label{eq:finite-place-odd-difference-map}
\end{equation}
and
\begin{equation}
 \mathcal Q_P=
 \frac{C(H_P^+,\mathbb Z)}
 {\sum_{p\in P}d_{p,P}(\mathcal F_{p,P})}.
 \label{eq:global-unit-fibre-cokernel-coordinate}
\end{equation}
There is a short exact sequence
\begin{equation}
 0\longrightarrow\mathcal Q_P
 \longrightarrow K_1(\mathfrak N_P)
 \longrightarrow\bigoplus_{p\in P}\mathcal C_{p,P}
 \longrightarrow0.
 \label{eq:natural-middle-K1-short-exact}
\end{equation}
No canonical splitting of \eqref{eq:natural-middle-K1-short-exact} is
asserted.
\end{theorem}

\begin{proof}
The quotient \(\mathfrak D_P\) is the finite direct sum of its boundary
strata.  Lemma~\ref{lem:unit-fibre-mapping-torus-k-groups} and
Theorem~\ref{thm:full-archimedean-unit-fibre-boundary} therefore give
\eqref{eq:global-natural-boundary-K0} and
\eqref{eq:global-natural-boundary-K1}.  Extension by zero from every edge to
the global one-zero algebra is a morphism of extensions.  Naturality and
additivity of the connecting maps, now applied to the complete edge groups
rather than only their units, combine
\eqref{eq:full-finite-edge-boundary-map} and
\eqref{eq:full-archimedean-boundary-map} into
\eqref{eq:global-full-unit-fibre-boundary-map}.

Since \(K_0(\mathfrak I_P)=0\) and
\(K_1(\mathfrak I_P)=C(K_P,\mathbb Z)\), the six-term sequence gives
\[
 K_0(\mathfrak N_P)=\ker\partial_P^{(1)}
\]
and
\begin{equation}
 0\to\operatorname{coker}\partial_P^{(1)}
 \to K_1(\mathfrak N_P)
 \to K_1(\mathfrak D_P)\to0.
 \label{eq:natural-six-term-middle-fragment}
\end{equation}
Equation \eqref{eq:global-full-unit-fibre-boundary-map} vanishes exactly when
\(S_h=-r_P^*h_\infty\).  The right side is sign-even.  Thus \(S_h\) must be
sign-even; when it is, the unique solution is
\(h_\infty=-S_h|_{H_P^+}\).  This proves
\eqref{eq:natural-middle-K0-exact-kernel} and its explicit inverse.

It remains to identify the cokernel without discarding the archimedean image.
Define
\begin{equation}
 \operatorname{Odd}_{H_P^+}:C(K_P,\mathbb Z)
 \longrightarrow C(H_P^+,\mathbb Z),
 \qquad
 \operatorname{Odd}_{H_P^+}(f)(c)=f(c)-f(-c).
 \label{eq:odd-difference-quotient-map}
\end{equation}
This map is onto: prescribe a function on \(H_P^+\) and extend it by zero on
\(-H_P^+\).  Its kernel is exactly the sign-even subgroup
\(r_P^*C(H_P^+,\mathbb Z)\).  Hence it induces the exact isomorphism
\[
 \frac{C(K_P,\mathbb Z)}{r_P^*C(H_P^+,\mathbb Z)}
 \xrightarrow{\ \cong\ }C(H_P^+,\mathbb Z).
\]
Under this map, the finite-place summand \(\iota_{p,P}(h)\) becomes exactly
\(d_{p,P}(h)\).  Therefore
\(\operatorname{coker}\partial_P^{(1)}\cong\mathcal Q_P\).  Substitution into
\eqref{eq:natural-six-term-middle-fragment} proves
\eqref{eq:natural-middle-K1-short-exact}.
\end{proof}

\paragraph{The one-prime coordinate check.}
If \(P=\{p\}\), then \(Y_{p,P}\), \(H_{p,P}\), and every finite quotient
are one point.  Thus
\(\mathcal F_{p,P}=\mathcal C_{p,P}=\mathbb Z\).  The finite edge has
\[
 K_0(A_{\varnothing p}^{(P)})=0,
 \qquad
 0\to C(K_P,\mathbb Z)/\mathbb Z\mathbf1_{K_P}
 \to K_1(A_{\varnothing p}^{(P)})\to\mathbb Z\to0.
\]
For the global natural algebra, the K-zero kernel consists exactly of
\((n,-n\mathbf1_{H_P^+})\), so \(K_0(\mathfrak N_P)\cong\mathbb Z\).  The
finite odd-difference map is zero, and therefore
\[
 0\to C(H_P^+,\mathbb Z)\to K_1(\mathfrak N_P)\to\mathbb Z\to0.
\]
This last sequence splits abstractly because its quotient is \(\mathbb Z\),
but the calculation selects no splitting.  The term
\(C(H_P^+,\mathbb Z)\) is the exact residual sign-asymmetric unit-fibre
coordinate after the archimedean boundary has filled every sign-even function.

\paragraph{Exact comparison with the stable balancing extension.}
For constant \(h_p=k_p\) and constant \(h_\infty=n\),
\eqref{eq:global-full-unit-fibre-boundary-map} reduces to
\[
 -\left(\sum_{p\in P}k_p+n\right)\mathbf1_{K_P}.
\]
The two stable archimedean character projections both map to the same
constant \(n=1\) class, exactly as proved in
Corollary~\ref{cor:stable-to-natural-K-square-lifted}.  Thus the earlier
codiagonal square is the constant-coordinate restriction of the full map,
not a substitute for it.  The new calculation also shows precisely what lies
outside that image: finite invariant functions, finite coinvariants with the
multiplicities \eqref{eq:finite-coinvariant-transition}, and the odd-difference
quotient \(\mathcal Q_P\).  It proves no surjectivity, Morita equivalence,
choice-independence across prime sets, endpoint statement, or global zeta
consequence.

\paragraph{Bounded finite-coordinate replay.}
The executable certificate and its receipt are stored under
\begin{center}
\footnotesize
\path{certificates/globalization_nte_nonconstant_unit_k/}.
\end{center}
The receipt has stable ID
\texttt{CERT-GNTE-NONCONSTANT-UNIT-K-20260830-0001}.
They verify the finite quotient orbit lengths, finite-level stabilizers, invariant
and coinvariant bases, both transition formulas including their
\(L_{\mathbf m'}/L_{\mathbf m}\) multiplicities, injectivity of the
transition matrices, the finite and archimedean pullback matrices, and the
global kernel/cokernel coordinate identities on bounded examples.  The
certificate does not replace the mapping-torus theorem, Green imprimitivity,
six-term exactness, or the proofs above, and it launches no Lean or Lake
process.
